% ============================================================================
% §3.2 (Screening ¶) + §3.3 ANSWER VERIFICATION — v3 draft (2026-07-24)
% Written in the post-07-23 register (technical, formal, direct; KDD-level).
%
% WHAT THIS DELIVERS
%   (1) §3.2 "Screening" paragraph, replacement. Now shows exactly what the
%       admissibility prover proves: the subterm S, the syntactic injectivity
%       argument, the surjectivity axiom as displayed first-order sentence,
%       the exact prover query, and the scope of the conclusion.
%       *** CORRECTNESS FIX (not just style): the compiled text says "that
%       subterm acts as a left inverse, making the map x -> T injective".
%       Both halves are off: the x-free CONTEXT C is the left inverse, and the
%       injective map is S (the subterm), not T. Fixed below; verify against
%       prove_status.py before merging. ***
%   (2) §3.3 full replacement, including the sentence at the literal
%       "Figure 3 ..." placeholder, the explicit three-step argument for the
%       construction channel, the "prover never reasons about infinity"
%       sentence, the self-policing paragraph, and the negative-controls
%       sentence (gate contents verified against llm_construct.py --selftest
%       on 2026-07-24: positive control certifies; left projection rejected at
%       check (a); collapsing presentation rejected at check (b); the parser
%       additionally rejects literal variable-to-variable equations and any
%       non-equational syntax).
%
% WHAT TO CHANGE BEFORE THE FIGURE-3 WALK-THROUGH (your question):
%   - PDF ¶ "A solved ALPS instance is a verified mathematical fact, not a
%     predicted label."  → replaced (second antithesis; abstract holds the
%     paper's one allocation). New opener also states the trust base, which
%     TASK_AND_JUDGING.md requires the paper to state and the compiled text
%     currently omits.
%   - PDF ¶ "The construction channel. ... In other words, E must impose
%     enough constraints..." → kept in substance, tightened; the semantic
%     reading of E now leads and the two claims are numbered so the Figure 3
%     caption and the formal queries can refer to them.
%   - The trivial-channel ¶ gains the header/footer/allowlist mechanism (in
%     your v2 draft, dropped from the compile) and one show-don't-tell
%     sentence giving the shape of a trivial-side proof (a chain of law
%     applications), which §4.2's assembler then picks up as a breadcrumb.
%
% Citation keys and labels follow section3_draft.tex: kisielewicz1997austin,
% baader1998term, moura2021lean, kovacs2013first, smallbone2021twee,
% eq:dichotomy, eq:trivialgoal, eq:construction. %% VERIFY the Figure 3 label
% (fig:toy below) against the Overleaf source.
% ============================================================================

% ---------------------------------------------------------------------------
% §3.2 — replacement for the "Screening." paragraph
% ---------------------------------------------------------------------------

\textbf{Screening.} Each candidate then passes screens ordered from cheap to
expensive (Figure~\ref{fig:funnel}). A fast proof search first removes
candidates it can prove trivial at generation time, and a finite-model filter
searches carriers of increasing size and discards every candidate it can
satisfy nontrivially. The survivors go to the admissibility prover, which
attempts a machine-checked proof that no nontrivial finite model exists. The
proof has a fixed shape. For a candidate $L\colon x = T[x,\bar y]$, the
screener selects a subterm $S$ of $T$ that contains every occurrence of $x$,
so that $T = C[S]$ for a context $C$ in which $x$ does not appear. The law
itself then reads $x = C[S(x,\bar y)]$: the context $C$ is a left inverse of
$S$, so the map $x \mapsto S(x,\bar y)$ is injective, with no proof search
required. On a finite carrier every injective map is surjective, so every
finite model of $L$ additionally satisfies the first-order sentence
\begin{equation}
\label{eq:surj}
\forall \bar y\ \forall u\ \exists x\colon\; S(x,\bar y) = u.
\end{equation}
The admissibility query submits $L$ together with \eqref{eq:surj} to the
prover and asks for $x = y$. A refutation proof establishes that every model
satisfying both is trivial; since every nontrivial finite model of $L$ would
satisfy both, none exists. The certificate is that refutation proof, stored
with the instance. Its scope is deliberately limited: \eqref{eq:surj} is
asserted only inside the finite-model argument, so the proof establishes
``no nontrivial finite model'' and never $L \models x = y$, leaving the
dichotomy of \eqref{eq:dichotomy} intact. Admission requires this
certificate, and no such certificate exists for a law with a nontrivial
finite model, so none can enter the corpus. The price is incompleteness:
when no subterm of $T$ contains every occurrence of $x$ and injectivity
cannot be proved by other means, the query certifies nothing, and the
candidate is discarded even though it may be admissible. What remains after
screening is the admissible pool, split by \eqref{eq:dichotomy} into laws
proved trivial, laws whose models the screening provers construct outright,
and the residual that the baseline of the next section attempts.

% ---------------------------------------------------------------------------
% §3.3 — full replacement
% ---------------------------------------------------------------------------

\subsection{Answer Verification}

Each solved ALPS instance is a machine-verified mathematical statement. The
two sides of the task call for different verification machinery, so the judge
issues a certificate through the channel that fits the answer: a proof
checked by the Lean kernel for the trivial side, and an automated-prover
certificate for a constructed model. Both channels are fully automatic, a
submission is accepted only when its check succeeds, and the trust base is
stated plainly: the Lean kernel on the trivial side
\citep{moura2021lean}, and prover saturation on the construction side
\citep{kovacs2013first}, cross-checkable by a second prover
\citep{smallbone2021twee} and the same base on which the corpus labels
themselves rest.

\textbf{The trivial channel.} For each law the judge generates the
proposition
\begin{equation}
\label{eq:trivialgoal}
\textsc{TrivialGoal} \equiv \forall M\ \forall {\op}\colon M \times M \to
M,\; \mathrm{Law}(\op) \rightarrow \forall a, b : M,\ a = b,
\end{equation}
where $\mathrm{Law}(\op)$ asserts that $\op$ satisfies $L$, and a submission
is a Lean proof of it. The statement is generated, never written by the
submitter: the judged file fixes \eqref{eq:trivialgoal} in a header, inserts
the submitted proof, and forces it to elaborate at exactly that type, and the
proof's axiom footprint must fall within a fixed allowlist. A proof the
kernel accepts under these constraints is correct by construction. In
practice a proof of \eqref{eq:trivialgoal} formalizes a chain of equalities
$a = t_1 = \cdots = t_k = b$ in which each step applies one instance of $L$;
the evaluation harness of Section~4.2 is built around this shape.

\textbf{The construction channel.} A solver that judges $L$ to be Austin
does not write a Lean proof about an infinite object. It submits the model
itself, as a finite presentation $E$: a set of equations over $\op$, distinct
from $L$, whose models are all the magmas that satisfy them. The submission
makes two claims about this class: (i) every model of $E$ satisfies $L$, and
(ii) some model of $E$ has two distinct elements. The two claims pull in
opposite directions — $E$ must impose enough structure to entail the law
without imposing so much that all elements collapse to one — and the
certificate requires both. Figure~\ref{fig:toy} traces the full certificate
on a toy law. Formally, the judge issues two prover queries
\citep{kovacs2013first},
\begin{equation}
\label{eq:construction}
\text{(a)}\;\; E \vdash L
\qquad\qquad
\text{(b)}\;\; E \cup \{\,a \neq b\,\} \text{ is satisfiable,}
\end{equation}
and certifies the submission only when both succeed. Check (a) returns a
refutation proof, which establishes claim (i). Check (b) is a saturation
that terminates without contradiction under a completeness-preserving
strategy; the saturated clause set describes a model of $E$ with two
distinct elements, which establishes claim (ii). The two checks compose
into a three-step argument: some magma satisfies $E$ with $a \neq b$, every
magma satisfying $E$ satisfies $L$, and therefore $L$ has a nontrivial
model, which the admissibility certificate of Section~3.2 forces to be
infinite. Neither query mentions infinity. Both are finite first-order
tasks, and infinitude enters only through the admissibility proof already
carried by the instance.

The pair of checks is resistant to degenerate submissions. Submitting the
law itself as $E$ gains nothing: check (a) becomes immediate, but check (b)
reduces to the bare saturation of $L$, which is the computation that fails
to terminate on the hard tier. A collapsing presentation such as
$\{x \op y = x,\ x \op y = y\}$ entails every law and passes check (a), but
fails check (b), because $E \cup \{a \neq b\}$ is unsatisfiable. A
presentation too weak to force the law, such as the left projection
$\{x \op y = x\}$, passes check (b) but fails check (a). The judge is
validated against controls of all three kinds before any evaluation run: a
known-correct presentation certifies, the left projection is rejected at
check (a), and the collapsing presentation is rejected at check (b). The
submission parser additionally rejects anything that is not a set of
unconditional equations, so implications, disequations, and literal
restatements of $x = y$ never reach the prover.

\textbf{Why two channels.} A single Lean judge covering both sides is not
available. If an admissible law had a model whose operation is a polynomial
or affine formula over $\mathbb{Z}$, $\mathbb{Z}^k$, or $\mathbb{Z}/n$, the
identity would survive reduction modulo two and yield a nontrivial
\emph{finite} model, contradicting admissibility. No Austin model is an
arithmetic formula, so the model families that are straightforward to
formalize in Lean are provably empty here, and an Austin answer cannot be a
small arithmetic expression recalled from training. The structures that do
satisfy these laws — term models built from saturations — await a
formalization of ground confluence, which we develop as companion work
\citep{janota2026case}. The two-channel judge is what makes the construction
side gradable today, at the cost of one additional trusted checker.

% ---------------------------------------------------------------------------
% NOTES (delete before merge)
% ---------------------------------------------------------------------------
% 1. eq:surj adds one numbered display to §3.2. If space is tight, inline it:
%    "the first-order sentence  ∀ȳ∀u∃x: S(x,ȳ)=u  (surjectivity of S)".
% 2. The collapsing-control example was changed from {x = y} to
%    {x◇y = x, x◇y = y}. The submission parser rejects the literal string
%    "x = y" (degenerate-equation gate in llm_construct.py), so {x = y} is a
%    control the parser, not the prover, would catch; the two-equation form
%    is what actually reaches check (b) and is rejected there. Figure 3's
%    caption uses the same logic on E0 and needs no change. If you prefer to
%    keep {x = y} for readability, add "(rejected at the parser in our
%    implementation)".
% 3. "distinct from L" in the construction ¶: strictly, the harness does not
%    forbid E = {L}; it is self-defeating rather than forbidden. Current
%    phrasing survives because the self-policing ¶ explains it, but if you
%    want exactness, replace "distinct from L" with "typically distinct
%    from L".
% 4. The trust-base sentence is required by TASK_AND_JUDGING.md ("state that
%    in the paper") and is currently absent from the compile.
% 5. Breadcrumb "the evaluation harness of Section 4.2 is built around this
%    shape" plants the assembler; keep §4.2's "the judge and assembler are
%    exactly the tools any solver built on ALPS would use" as the payoff.
% 6. Terminology check against the fixed list: "checks (a) and (b)",
%    "presentation E", "admissible", "hard tier" all used; no "horns".
% ============================================================================
